\section*{Critical slack as an influence-weighted child deficit}

Let $f:\{-1,1\}^n\to\{-1,1\}$ be nonconstant and write
$R=\mathsf H(f)/\mathsf I(f)$.  For a relevant coordinate $i$, let
$f_i^\sigma$ be the two physical children, and put
\[
 H_\sigma=\mathsf H(f_i^\sigma),\quad
 I_\sigma=\mathsf I(f_i^\sigma),\quad
 r_\sigma=H_\sigma/I_\sigma
\]
when $I_\sigma>0$.  Set
\[
 \overline H_i={H_++H_-\over2},\qquad
 \overline I_i={I_++I_-\over2},\qquad
 \Delta_iH=\mathsf H(f)-\overline H_i.
\]

\begin{theorem}[critical-slack decomposition]
For every relevant coordinate,
\[
 \boxed{
 \Delta_iH-RI_i(f)
 ={1\over2}\sum_{\sigma:I_\sigma>0}
 I_\sigma\bigl(R-r_\sigma\bigr).}                 \tag{1}
\]
If $f$ is descendant-critical, then the left side is nonnegative.  Moreover,
for every $t>0$,
\[
 {1\over2}\sum_{\sigma:r_\sigma\le R-t}I_\sigma
 \le {\Delta_iH-RI_i(f)\over t}.                 \tag{2}
\]
Consequently zero slack at coordinate $i$ forces every nonconstant child to
have ratio exactly $R$.  More generally, if
$\Delta_iH-RI_i(f)\le\varepsilon\overline I_i$, then under the child law
weighted by child influence, the expected ratio loss is at most
$\varepsilon$ and the probability of loss at least $t$ is at most
$\varepsilon/t$.
\end{theorem}

\begin{proof}
Physical edge counting gives
\[
 \overline I_i=\mathsf I(f)-I_i(f).
\]
Using $\mathsf H(f)=R\mathsf I(f)$,
\begin{align*}
 \Delta_iH-RI_i(f)
 &=\mathsf H(f)-\overline H_i-RI_i(f)\\
 &=R\overline I_i-\overline H_i\\
 &={1\over2}\sum_\sigma(RI_\sigma-H_\sigma),
\end{align*}
which is (1), with constant children contributing zero.  Descendant
criticality gives $r_\sigma\le R$ for each nonconstant child, so all terms
are nonnegative.  On the subset with $R-r_\sigma\ge t$, each summand is at
least $tI_\sigma$; this proves (2).  The remaining statements are the
equality case and Markov's inequality under normalization by
$2\overline I_i=I_++I_-$.
\end{proof}

\paragraph{Tensor compatibility.}
If $F=fg$ on disjoint blocks and $i$ belongs to the $f$ block, then
$I_i(F)=I_i(f)$ and $\Delta_iH(F)=\Delta_iH(f)$.  If $f,g$ are critical
with common ratio $R$, the entire slack identity for $F$ at $i$ equals that
for $f$.  Hence the theorem descends exactly to the equal-ratio factors from
the tensor-purity theorem; it does not pretend that tensor replication makes
the gap easier.

\paragraph{Scope.}
Abstract Fourier-support components are not physical children and cannot be
inserted into (1).  A one-point-surgery lift need not be critical; after
passing to a maximum-ratio physical descendant, (1) applies to the selected
descendant, not automatically to the unreduced surgery.  No universal upper
bound on the slack is asserted.
